\documentclass[11pt]{article}
\usepackage[T1]{fontenc}
\usepackage[utf8]{inputenc}
\usepackage{lmodern}
\usepackage{amsmath,amssymb}
\usepackage[margin=1in]{geometry}
\usepackage{booktabs}
\usepackage{float}
\usepackage{xcolor}
\usepackage[hidelinks]{hyperref}
\definecolor{baseline}{HTML}{24557A}
\definecolor{sketch}{HTML}{A3442B}
\setlength{\parskip}{0.35em}
\setlength{\parindent}{0pt}
\emergencystretch=2em

% The vector plot below places 25-origin bin means from the archived
% pilot_results/trajectory.csv on a 360 by 180 point plotting area.

\title{Volatility-Estimation Lag and Calibration Memory:\\
A Sequential Prediction-Interval Feasibility Study}
\author{Junguang Jia}
\date{}

\begin{document}
\maketitle

\section*{Abstract}
Sequential prediction intervals can respond slowly to a scale change for two
reasons: the current scale forecast may lag, and earlier scale errors remain in
the calibration scores after that forecast recovers. This report separates
those mechanisms in a finite-window model and examines one controlled Gaussian
scale-jump pilot. The archived pilot contains 200 independent paths. For a
250-score window, the causal EWMA-normalized rolling interval attained 95.24\%
coverage overall but 83.42\% over the first 25 forecasts after the upward
scale step; the raw rolling interval attained 94.28\% and 53.72\%,
respectively. A separate deterministic contamination calculation shows how
inflated historical scores can persist after the current scale becomes exact.
These are development results from one synthetic design. The work does not
establish a new method, a general conditional-coverage guarantee, or a
novelty claim.

\section{Introduction}
This study concerns one-step intervals for a scalar next response $Y_t$ under
time-varying conditional scale. The question is whether scale normalization
improves the immediate response to a volatility change and how long historical
scale-estimation errors can affect a finite calibration window. The distinction
matters because an interval can have near-nominal coverage averaged over a long
evaluation period while missing many responses immediately after a shift.

Adaptive conformal inference (ACI) already studies online coverage under
distribution shift and uses GARCH-normalized scores in a realized-volatility
example~\cite{gibbs2021}. Sequential predictive conformal inference (SPCI)
predicts future residual quantiles using temporal dependence and gives an
assumption-dependent asymptotic conditional-coverage analysis~\cite{xu2023}.
Strongly adaptive online conformal prediction and conformal PID control provide
further adaptive calibration approaches~\cite{bhatnagar2023,angelopoulos2023}.
The present pilot implements none of those comparators. Its purpose is to
isolate a candidate mechanism and specify the evidence needed for a later
comparison, rather than to claim priority or superiority.

\section{Sequential interval protocol}
Let $\mathcal F_{t-1}$ contain the observations available before $Y_t$ is
revealed. At origin $t$, a deployable method uses only $\mathcal F_{t-1}$ to
issue a location forecast $\widehat\mu_t$, a positive scale forecast
$\widehat\sigma_t$, and an interval. After $Y_t$ is observed, it stores the
\emph{vintage} score
\[
 R_t=\frac{|Y_t-\widehat\mu_t|}{\widehat\sigma_t}.
\]
The forecasts in an old score remain those issued at its original origin;
they are not recomputed using a later model fit. In the pilot the mean is
known to be zero. For a rolling window of $W$ previous scores, let
\[
 k=\left\lceil(W+1)(1-\alpha)\right\rceil,\qquad
 q_t=R_{(k);\,t-W:t-1},\qquad
 C_t=[\widehat\mu_t-\widehat\sigma_t q_t,
      \widehat\mu_t+\widehat\sigma_t q_t].
\]
Here $R_{(k);\,t-W:t-1}$ is the $k$th order statistic of the saved scores.
If $k>W$, the specified threshold is infinite. The current response never
enters its own calibration window. Forecast cutoffs, unavailable scores, and
failed fits must remain visible in a later full evaluation; the corrected
rank alone gives no general coverage guarantee under temporal dependence.

\section{Scale error and calibration memory}
Suppose $Y_t=\mu_t+\sigma_t Z_t$ with exact location forecasts. If every
forecast scale in a window, including the current one, equals the true scale
times the same constant $c>0$, then $R_i=|Z_i|/c$. The score quantile is
$q_t=q_t^{\rm oracle}/c$, and
$\widehat\sigma_tq_t=\sigma_tq_t^{\rm oracle}$.
Thus a stable multiplicative scale bias cancels from this rolling interval.
Scale forecast error by itself does not determine interval error; the
\emph{change} in scale error across score vintages is relevant.

Write $\eta_i=\log(\widehat\sigma_i/\sigma_i)$ and let $G$ be the CDF of
$|Z_i|$. Assume the innovations are identically distributed and continuous,
each $Z_i$ is independent of $\mathcal F_{i-1}$, and each $\eta_i$ is fixed
before $Y_i$ is revealed. Define
$H_i(x)=G(e^{\eta_i}x)$, its window average $\overline H_t(x)$, and the
empirical score CDF $\widehat H_t(x)$. When the $W$ scores are distinct and
$k\leq W$, the conditional coverage $c_t=H_t(q_t)$ obeys the pathwise
decomposition
\[
 |c_t-(1-\alpha)|\leq
 \underbrace{\sup_x|\widehat H_t(x)-\overline H_t(x)|}_{A_t}
 +\underbrace{\sup_x|\overline H_t(x)-H_t(x)|}_{B_t}
 +\left|\frac{k}{W}-(1-\alpha)\right|.
\]
If $G$ has density $g$ and $L=\sup_{u\geq0}ug(u)<\infty$, then
$B_t\leq L W^{-1}\sum_{i=t-W}^{t-1}|\eta_i-\eta_t|$.
This separates empirical fluctuation, relative scale-error drift, and rank
rounding. It is a diagnostic decomposition, not an automatic finite-sample
conditional-coverage theorem; the empirical term requires its own analysis.

For an exact, narrower calculation, assume the current scale has recovered.
Among $W$ independent historical scores, $W-m$ have distribution $U=|Z|$ and
$m$ have distribution $cU$. Membership in the two groups and $c$ are fixed
before the scores are drawn. Let $q$ be their $k$th order statistic and let
$U_0$ be an independent new draw. Conditioning on $U_0=u$ gives
\begin{equation}
 \Pr(U_0\leq q)=\int_0^\infty
 \Pr\{A_u+B_u\leq k-1\}\,dG(u),\qquad
 \begin{cases}
 A_u\sim\operatorname{Bin}(W-m,G(u)),\\
 B_u\sim\operatorname{Bin}(m,G(u/c)),
 \end{cases}
 \label{eq:contamination}
\end{equation}
with independent binomial counts for each fixed $u$. Indeed, $U_0\leq q$
exactly when at most $k-1$ historical scores lie below $U_0$, apart from
probability-zero ties. For $m=0$ or $c=1$, the clean repeated-sample coverage
is $k/(W+1)$. Coupling the same $U$ draws shows that $c>1$ cannot lower this
repeated-sample coverage, while $c<1$ cannot raise it. This identity
does not transfer unchanged to adaptive EWMA or fitted GARCH scores, whose
errors depend on past observations. Its proof and relation to prior work still
require independent scientific review.

\section{Exploratory experimental design}
The archived pilot draws 200 independent paths of length 2,600 with
$Z_t\sim N(0,1)$, known mean zero, and true scale $\sigma_t=1$ until
zero-based origin 1200, $3$ for origins 1200--1799, then $1$ again. Evaluation
covers origins 1000--2599 at nominal coverage $1-\alpha=0.95$. The causal
EWMA forecast starts at variance $v_0=1$ and updates only after reveal:
\[
 \widehat\sigma_t=\sqrt{v_t},\qquad
 v_{t+1}=0.94v_t+0.06Y_t^2.
\]
The principal comparison uses raw rolling absolute errors, EWMA-normalized
rolling vintage scores, and an oracle-scale rolling diagnostic, all with
$W=250$. A full Gaussian oracle uses the true current scale and the normal
0.975 quantile. All methods share the same innovations within each path;
oracle rows have privileged information and are not deployable competitors.

Coverage is averaged first across forecast origins within each path and then
across paths. Its Monte Carlo standard error (MCSE) is the sample standard
deviation of the 200 path-level means divided by $\sqrt{200}$. Mean width is
reported relative to the full Gaussian oracle width at the same origin.
Analytic conditional coverage in this Gaussian simulation is also averaged
over paths for the event-time figure. These uncertainty summaries condition
on this single fixed data-generating design, not a population of
volatility processes. The archive also contains $W=50$ variants; their clean
rank benchmark is $49/51=0.9608$, versus $239/251=0.9522$ for $W=250$,
so their coverage difference cannot be interpreted as pure adaptivity.

\section{Results}
Table~\ref{tab:pilot} shows that average coverage can mask a severe immediate
post-shift deficit. For raw rolling calibration, the first 25 upward-shift
forecasts cover 53.72\% of outcomes; EWMA normalization raises this to
83.42\% in this design but still falls well below 95\%. The oracle-scale
diagnostic is close to its clean finite-rank level. Following the downward
step, raw rolling coverage over the first 25 forecasts is 100\%, with mean
width 3.03 times the full-oracle width; the EWMA-normalized counterpart has
99.98\% coverage and width ratio 2.38. Neither observation establishes a
method ranking against strong adaptive baselines.

\begin{table}[H]
\centering
\small
\caption{Archived 200-path pilot at $W=250$. Coverage and MCSE are percentage
points; width ratio is the mean interval width divided by full Gaussian oracle
width. Both oracles use the true current scale.}
\label{tab:pilot}
\begin{tabular}{lrrrrr}
\toprule
Method & Overall (\%) & MCSE & Up 0--24 (\%) & MCSE & Width ratio \\
\midrule
Raw rolling & 94.283 & 0.027 & 53.72 & 0.570 & 1.195 \\
EWMA-normalized rolling & 95.242 & 0.015 & 83.42 & 0.355 & 1.079 \\
Oracle-scale rolling & 95.225 & 0.017 & 95.40 & 0.282 & 1.017 \\
Full Gaussian oracle & 94.962 & 0.041 & 94.90 & 0.301 & 1.000 \\
\bottomrule
\end{tabular}
\end{table}

Figure~\ref{fig:upshift} shows a second view of the same upward shift using
conditional coverage, which averages over the future Gaussian outcome while
retaining variation in past histories. The plotted values are descriptive
25-origin bin means from the archived trajectory. They do not supply
bin-level uncertainty intervals or a new confirmatory test.

\begin{figure}[H]
\centering
\setlength{\unitlength}{1pt}
\begin{picture}(440,250)
{\color{gray!18}
\multiput(40,30)(0,36){6}{\line(1,0){360}}
\multiput(40,30)(72,0){6}{\line(0,1){180}}}
\put(40,30){\line(1,0){360}}
\put(40,30){\line(0,1){190}}
\put(40,12){\makebox(0,0){\footnotesize 0}}
\put(112,12){\makebox(0,0){\footnotesize 50}}
\put(184,12){\makebox(0,0){\footnotesize 100}}
\put(256,12){\makebox(0,0){\footnotesize 150}}
\put(328,12){\makebox(0,0){\footnotesize 200}}
\put(400,12){\makebox(0,0){\footnotesize 250}}
\put(32,30){\makebox(0,0)[r]{\footnotesize 0.5}}
\put(32,66){\makebox(0,0)[r]{\footnotesize 0.6}}
\put(32,102){\makebox(0,0)[r]{\footnotesize 0.7}}
\put(32,138){\makebox(0,0)[r]{\footnotesize 0.8}}
\put(32,174){\makebox(0,0)[r]{\footnotesize 0.9}}
\put(32,210){\makebox(0,0)[r]{\footnotesize 1.0}}
\put(40,233){\makebox(0,0)[l]{\small Mean conditional coverage}}
\put(220,0){\makebox(0,0){\small Forecast origins after upward scale step}}
{\color{baseline}\linethickness{0.8pt}
\qbezier(58,49)(76,76.5)(94,104)
\qbezier(94,104)(112,123)(130,142)
\qbezier(130,142)(148,151.5)(166,161)
\qbezier(166,161)(184,166.5)(202,172)
\qbezier(202,172)(220,175.5)(238,179)
\qbezier(238,179)(256,181.5)(274,184)
\qbezier(274,184)(292,185.5)(310,187)
\qbezier(310,187)(328,188.5)(346,190)
\qbezier(346,190)(364,191)(382,192)
\put(58,49){\circle*{3}}\put(94,104){\circle*{3}}
\put(130,142){\circle*{3}}\put(166,161){\circle*{3}}
\put(202,172){\circle*{3}}\put(238,179){\circle*{3}}
\put(274,184){\circle*{3}}\put(310,187){\circle*{3}}
\put(346,190){\circle*{3}}\put(382,192){\circle*{3}}
\put(265,108){\line(1,0){17}}\put(273,108){\circle*{3}}}
\put(287,108){\makebox(0,0)[l]{\footnotesize Raw rolling}}
{\color{sketch}\linethickness{0.8pt}
\qbezier(58,153)(76,174)(94,195)
\qbezier(94,195)(112,196)(130,197)
\qbezier(130,197)(148,197.5)(166,198)
\qbezier(166,198)(184,198)(202,198)
\qbezier(202,198)(220,198)(238,198)
\qbezier(238,198)(256,198)(274,198)
\qbezier(274,198)(292,198)(310,198)
\qbezier(310,198)(328,197.5)(346,197)
\qbezier(346,197)(364,197)(382,197)
\put(58,153){\makebox(0,0){\scriptsize$\diamond$}}
\put(94,195){\makebox(0,0){\scriptsize$\diamond$}}
\put(130,197){\makebox(0,0){\scriptsize$\diamond$}}
\put(166,198){\makebox(0,0){\scriptsize$\diamond$}}
\put(202,198){\makebox(0,0){\scriptsize$\diamond$}}
\put(238,198){\makebox(0,0){\scriptsize$\diamond$}}
\put(274,198){\makebox(0,0){\scriptsize$\diamond$}}
\put(310,198){\makebox(0,0){\scriptsize$\diamond$}}
\put(346,197){\makebox(0,0){\scriptsize$\diamond$}}
\put(382,197){\makebox(0,0){\scriptsize$\diamond$}}
\put(265,91){\line(1,0){17}}\put(273,91){\makebox(0,0){\scriptsize$\diamond$}}}
\put(287,91){\makebox(0,0)[l]{\footnotesize EWMA-normalized}}
{\color{gray!70!black}
\multiput(58,193)(9,0){36}{\line(1,0){5}}
\multiput(265,74)(8,0){2}{\line(1,0){5}}}
\put(287,74){\makebox(0,0)[l]{\footnotesize Oracle-scale rolling}}
\end{picture}
\caption{Mean conditional coverage in consecutive 25-origin bins after the
upward scale step, averaged over the 200 archived paths. Points are placed at
bin midpoints. Lines guide the eye; no bin-level MCSE is shown.}
\label{fig:upshift}
\end{figure}

Table~\ref{tab:contamination} evaluates the independent deterministic-score
model in Eq.~\eqref{eq:contamination}, not the EWMA pilot. With $W=250$ and
$k=239$, 25 inflated historical scores at $c=3$ raise repeated-sample
coverage from the clean 95.219\% to 98.955\% despite an exact current scale.
Replacing $c$ by $1/3$ lowers it to 94.690\%. The numerical quadrature and
40,000 independent Monte Carlo checks agree within the recorded Monte Carlo
errors. These calculations illustrate a possible post-recovery memory effect;
they do not measure how often a fitted scale model creates that pattern.

\begin{table}[H]
\centering
\small
\caption{Deterministic contamination model at $W=250$, $k=239$. Quadrature
and Monte Carlo (MC) coverage are percentages; MCSE is in percentage points.}
\label{tab:contamination}
\begin{tabular}{rrrrr}
\toprule
Contaminated $m$ & Multiplier $c$ & Quadrature (\%) & MC (\%) & MCSE \\
\midrule
0 & 1 & 95.2191 & 95.2182 & 0.0067 \\
25 & 3 & 98.9553 & 98.9564 & 0.0037 \\
25 & $1/3$ & 94.6903 & 94.6942 & 0.0074 \\
50 & 3 & 99.9192 & 99.9187 & 0.0007 \\
\bottomrule
\end{tabular}
\end{table}

\section{Limitations and next study}
The pilot has one Gaussian scale-jump profile, one EWMA decay, a known mean,
and development seeds only. It includes neither a fitted GARCH scale forecast
nor ACI, SPCI, strongly adaptive online conformal prediction, or conformal
PID. It uses no dated real series. Oracle methods know the true current scale
at an unannounced change, so their performance is a diagnostic boundary, not
an achievable forecast at that origin. Overall coverage, interval width, and
event-time behavior must be evaluated together; no result here establishes
general conditional coverage or practical superiority.

The next study should first test Eq.~\eqref{eq:contamination} over a
predeclared grid of contamination counts and multipliers, including both
directions and clean controls. It should then add causally fitted scale
forecasts and faithful adaptive comparators on common planned evaluation
origins, retaining failed fits and infinite intervals. Hyperparameters,
scenarios, primary contrasts, and confirmation seeds should be fixed before
the confirmation runs. Mathematical assumptions and overlap with prior work
need independent review before a new contribution is claimed.

\section{Reproducibility}
The repository includes the Python pilot, numerical identity checks,
forecast-before-reveal tests, a locked dependency file, and the archived
synthetic outputs. The numerical source for Table~\ref{tab:pilot} is
\path{reference/synthetic_baseline_v1/pilot_results/summary.csv};
Figure~\ref{fig:upshift} aggregates
\path{reference/synthetic_baseline_v1/pilot_results/trajectory.csv}; and
Table~\ref{tab:contamination} uses
\path{reference/synthetic_baseline_v1/theory_results/coverage_identity.csv}.
The six numerical archive files have recorded SHA-256 hashes in
\path{reference/synthetic_baseline_v1/SHA256SUMS.txt}. This standalone source
embeds the plotted values and requires no external data file to compile.

The archived 200-path pilot records Python 3.13.5, NumPy 2.3.5, SciPy 1.17.0,
and seed 20260928. The current locked workspace uses Python 3.12.9. Its 56
Python tests passed; a fresh two-path development smoke run matched all 560
keyed per-replication metrics in the archived seed prefix at absolute and
relative tolerance $10^{-12}$. A full 200-path reproduction in the current
workspace has not been completed. The two-path agreement checks implementation
consistency for that prefix, not the research claims in this report.

\begin{thebibliography}{9}
\bibitem{gibbs2021}
I. Gibbs and E. J. Cand\`es.
Adaptive conformal inference under distribution shift.
\emph{Advances in Neural Information Processing Systems}, 34, 2021.
\href{https://arxiv.org/abs/2106.00170}{arXiv:2106.00170}.

\bibitem{xu2023}
C. Xu and Y. Xie.
Sequential predictive conformal inference for time series.
\emph{Proceedings of the 40th International Conference on Machine Learning},
PMLR 202:38707--38727, 2023.
\href{https://proceedings.mlr.press/v202/xu23r.html}{PMLR: xu23r}.

\bibitem{bhatnagar2023}
A. Bhatnagar, H. Wang, C. Xiong, and Y. Bai.
Improved online conformal prediction via strongly adaptive online learning.
\emph{Proceedings of the 40th International Conference on Machine Learning},
PMLR 202:2337--2363, 2023.
\href{https://proceedings.mlr.press/v202/bhatnagar23a.html}{PMLR: bhatnagar23a}.

\bibitem{angelopoulos2023}
A. N. Angelopoulos, E. J. Cand\`es, and R. J. Tibshirani.
Conformal PID control for time series prediction.
\emph{Advances in Neural Information Processing Systems}, 36, 2023.
\href{https://arxiv.org/abs/2307.16895}{arXiv:2307.16895}.
\end{thebibliography}
\end{document}
